% =====================================================================
%  tr09a.tex —— 原文 pp.114--126（译者：tr-04）
%  内容：§20 第五讲（+ 构造与一般 K 群的定义）；§21 第六讲（若干定理与例子）
%  原文无编号公式：展示式一律 \[ … \]；定理类号数照原文手工给出；
%  交换图用 tikzcd（照页图重排）。
% =====================================================================
% ===== 原文 pp.114--126 =====
% [原文 p.114]
\section{第五讲（+ 构造与一般 K 群的定义）}

我们用「空间」一词指带基点的 CW 复形，并用 $[X,Y]$ 表示保基点映射 $X\to Y$ 的同伦类集合。

设 $X$ 是一个空间，其基本群 $\pi_1(X)$ 是完美的，即 $\pi_1(X)_{ab}=0=H_1(X;\Z)$。

\para{问题：} 我们希望构造一个空间 $X^+$，使得 $\pi_1(X^+)=0$，同时有一个映射 $i:X\to X^+$，使诱导映射 $i_*:H_*(X)\to H_*(X^+)$ 是同调中的\emph{同构}。

\para{第一步.} 我们选取元素 $\gamma_i\in\pi_1(X)$（$i\in I$），使得它们生成的正规子群是整个群 $\pi_1(X)$。

我们假定如下特殊情形：某个 $\gamma$ 由环路 $S^1\xra{u}X$ 实现，并令
\[ Y=X\cup D^2, \]
其中 $D^2$ 是一个 2-圆盘（即一个 $2$-胞腔），其边界是环路 $\gamma$。

由 Van Kampen 定理（见附录）可得
\[ \pi_1(Y)=\pi_1(X)*\pi_1(D^2)=\pi_1(X)/\Gamma=0, \]
因为 $\pi_1(D^2)=0$，且其中 $\Gamma$ 表示由 $\gamma$ 生成的正规子群。于是：
\[ H_2(X)\to H_2(Y)\to H_2(X,Y)\to H_1(X)\to H_1(Y)\to 0. \]
但 $H_2(X,Y)=\Z$ 而 $H_1(X)=0$。因此
\[ H_n(X)\simeq H_n(Y),\quad n\ge 3. \]
有正合列
\[ 0\to H_2(X)\to H_2(Y)\xra{p}\Z\to 0. \]

% [Tutor Note: p.114 $\pi_1(X)_{ab}=0=H_1(X;\Z)$：第一个等号是 Hurewicz 定理（$\pi_1$ 的交换化同构于 $H_1$），所以这里「完美」与 $H_1=0$ 是同义语。+ 构造只对完美群（更一般地，完美正规子群）才可能保持同调，这正是本讲开头的假设。]

% [Tutor Note: p.114 原文只详述了「一个 $\gamma$」的特殊情形；一般情形（第一步中取一组 $\gamma_i$）要对每个 $\gamma_i$ 同时粘一个 2-胞腔，再对所有新出现的 $\pi_2$ 类粘 3-胞腔，并保证每一步都不改变已有同调——这就是 + 构造的全部内容，见 p.118 定理 1。]

% [原文 p.115]
我们应用 \emph{Hurewicz} 定理（见附录）得到：$\pi_1(Y)=0\Rightarrow\pi_2(Y)=H_2(Y)$。存在映射
\[ v:S^2\to Y \]
使得
\[ H_2(S^2)\to H_2(Y)\xra{p}\Z. \]
令
\[ X^+:=Y\cup_v C^3, \]
其中 $C^3$ 是一个 $3$-胞腔。于是得到如下序列：
\[ H_3(X^+)\to H_3(X^+,Y)\xra{\delta}H_2(Y)\to H_2(X^+)\to H_2(X^+,Y), \]
其中
\[ \delta:H_3(X^+,Y)\to H_2(Y), \]
而 $\delta$ 的像就是 $v$ 的类。因此
\[ H_2(X)\to H_2(X^+) \]
且
\[ H_n(X)=H_n(Y)=H_n(X^+),\quad n\ge 3. \]

\begin{prop}{1} 假定 $\pi_1(X)$ 是完美的，设 $i:X\to X^+$ 使得 $\pi_1(X^+)=0$，且诱导映射 $i_*:H_*(X)\to H_*(X^+)$ 是同构。那么对任意满足 $\pi_1(Y)=0$ 的 $Y$，
\[ i^*:[X^+,Y]\to[X,Y] \]
也是同构。 \end{prop}

\begin{cor}{1} $(X^+,i)$ 只在同伦等价的意义下确定。 \end{cor}

\begin{proof}[命题 1 的证明]
为证明满射性，考虑如下图表：
\[
\begin{tikzcd}
X \arrow[r,"i"] \arrow[d,"u"] & X^+ \arrow[d,"u'"] \\
Y \arrow[r,"i'"] & {Y\cup_X X^+=Z}
\end{tikzcd}
\]

% [Tutor Note: p.115 序列 $H_3(X^+)\to H_3(X^+,Y)\xra{\delta}H_2(Y)\to H_2(X^+)\to\cdots$ 来自对 $(X^+,Y)$ 的整系数同调长正合列；$\delta$ 把 3-胞腔的基本类送到环路 $\gamma$ 之实现 $v$ 的类，故 $H_2(X^+)\cong H_2(Y)/\langle[v]\rangle$，而 $H_2(X)\to H_2(X^+)$ 说明同调未被改变。]

% [原文 p.116]
若 $i$ 是同调同构，则 $i'$ 也是。

由 Van Kampen 定理有 $\pi_1(Z)=0$。

由 Whitehead 定理（单连通版本，见附录）得到 $i':Y\to Z$ 是同伦等价。现在令 $\tau:Z\to Y$ 为 $i'$ 的同伦逆。则
\[ (\tau u')i=\tau(i'u)=(\tau i')u\sim(\id)u=u. \]
这证明了 $i^*:[X^+,Y]\to[X,Y]$ 是满射。

为证明单射性，考虑如下图表：
\[
\begin{tikzcd}
X \arrow[r,"i"] \arrow[d,"u"] & X^+ \arrow[d,"u'"] \\
Y \arrow[r,"i'"] & Z
\end{tikzcd}
\]

接下来考虑两个映射
\[ g_0,g_1:X^+\to Y, \]
使得 $g_0i$ 与 $g_1i$ 都同伦于 $u$。然后我们应用同伦扩张定理，把 $g_0$ 与 $g_1$ 形变到 $g_0i=g_1i=u$。于是定义
\[ \tau_0:Z\xra{(\id_Y,g_0)}Y \]
以及
\[ \tau_1:Z\xra{(\id_Y,g_1)}Y \]
使得
\[ \tau_0i'=\tau_1i'=\id_Y. \]
但 $i'$ 是同伦等价，故 $\tau_0\sim\tau_1$。由于 $\tau_0u'=g_0$ 且 $\tau_1u'=g_1$，我们推出 $g_0\sim g_1$。

% [Tutor Note: p.116 证明满射性时用了同伦扩张定理：$g_0,g_1$ 在 $X$ 上（经 $i$）的限制都已同伦于 $u$，故可先形变使 $g_0i=g_1i=u$ 再作论证；原文只以 "we deform" 一句带过。Whitehead 定理此处需单连通版本，$Y,Z$ 单连通正是为此。]

% [原文 p.117]
\end{proof}

假设 $G$ 是完美的，并令
\[ 1\to H_2(G)\to\wt G\to G\to 1 \]
为其万有中心扩张。回忆 $H_1(G)=H_2(\wt G)=0$。于是有下述结果：

\begin{lem}{1}（令 $C=H_2(G)$。）存在纤维化之间的一个映射：
\[
\begin{tikzcd}
BC \arrow[r] \arrow[d,equal] & B\wt G \arrow[r] \arrow[d] & BG \arrow[d] \\
BC \arrow[r] & B\wt G^+ \arrow[r] & BG^+
\end{tikzcd}
\]
\end{lem}

\begin{proof}
（回忆 $\pi_1(BG)=G$。）证明由图表追踪给出：
\[ H^2(G,C)=H^2(BG,C)=[BG,K^{EM}(C,2)] \]
（其中 $K^{EM}$ 是 Eilenberg--McLane 空间）。此外
\[
\begin{tikzcd}
BC \arrow[r,equal] \arrow[d] & BC \arrow[r,equal] \arrow[d] & {K^{EM}(C,1)} \arrow[d] \\
B\wt G \arrow[r,"\beta"] \arrow[d] & P \arrow[r] \arrow[d] & {E(C,2)} \arrow[d] \\
BG \arrow[r] & BG^+ \arrow[r,"\alpha"'] & {K^{EM}(C,2)}
\end{tikzcd}
\]
其中 $\alpha:BG^+\to K^{EM}(C,2)$ 由前面的命题 1 存在，$\beta:B\wt G\to P$。我们还有：
\[ \pi_2(BG^+)=H_2(BG^+)=H_2(BG)=C. \]
\end{proof}

% [Tutor Note: p.117 原文把 Eilenberg--MacLane 拼作 "Eilenberg-McLane"（标准拼法为 MacLane），正文照原样；$K^{EM}$ 只在这一段使用。图中两处双线 "=" 是恒等映射（tikzcd 的 equal 箭头）；$B\wt G^+$ 指对空间 $B\wt G$ 作 + 构造。]

% [原文 p.118]
但 $\alpha$ 在 $\pi_2$ 上诱导同构：
\[ \pi_2(BG^+)\to\pi_1(BC)\to\pi_1(P)\to\pi_1(BG^+), \]
其中 $\pi_1(BC)=C$ 且 $\pi_1(BG^+)=0$。因此 $\pi_1(P)=0$。

现在 $\beta$ 是同调同构，$\pi_1(P)=0\Rightarrow B\wt G^+\sim P$，即二者同伦等价。

由映射
\[ BC\to B\wt G^+\to BG^+, \]
利用 Hurewicz 定理，我们发现
\[ \pi_3(BG^+)=\pi_3(B\wt G^+)=H_3(B\wt G^+), \]
因为 $H_1(\wt G)=H_2(\wt G)=0$。

下面假设 $G=E(A)$ 是完美的。我们有 $\wt G=\mathrm{St}(A)$ 以及
\[ H_2(G)=K_2(A). \]

\para{断言：}
\[ \pi_2(BE(A)^+)=K_2(A)=H_2(E(A)), \]
\[ \pi_3(BE(A)^+)=H_3(\mathrm{St}(A)). \]

\begin{proof}
该断言的证明由上面提到的内容推出：
\[ \pi_2(BE(A)^+)=H_2(BE(A)^+)=H_2(BE(A))=H_2(E(A))=K_2(A), \]
\[ \pi_3(BE(A)^+)=\pi_3(B\mathrm{St}(A)^+)=H_3(B\mathrm{St}(A))=H_3(\mathrm{St}(A)). \]
\end{proof}

\begin{thm}{1} 设 $N\subset\pi_1(X)$ 是完美正规子群。则存在一个空间 $X^+$（依赖于 $N$）以及一个映射 $i:X\to X^+$，使得：
\begin{itemize}[leftmargin=1.6em,itemsep=2pt]
\item 群 $\pi_1(X)/N\simeq\pi_1(X^+)$ 同构。
% [原文 p.119]
\item 对任意 $\pi_1(X^+)$-模 $L$，同调中的诱导映射
\[ i_*:H_*(X,i^*L)\to H_*(X,L) \]
% [Tutor Note: p.119 该式右端原文作 $H_*(X,L)$，但按 $i:X\to X^+$ 与「$L$ 是 $\pi_1(X^+)$-模、$i^*L$ 为其沿 $i$ 的拉回」的设定，右端应为 $H_*(X^+,L)$（对照同页 $H_*(B\GL A,L)\simeq H_*(B\GL A^+,L)$）；已用 800 dpi 裁图复核页图确印作 $H_*(X,L)$，故正文照原样，此处给出正确形式。]
是同构。
\end{itemize}
\end{thm}

二元组 $(X^+,i)$ 只在同伦等价的意义下确定。

\begin{proof}
本质上我们将构造出具有所需性质的空间 $X^+$，这就是 \emph{Quillen 著名的 + 构造}。

设 $\wt X$ 是对应于子群 $N\subset\pi_1(X)$ 的覆盖空间，即 $\pi_1(\wt X)=N$。于是有：
\[
\begin{tikzcd}
\wt X \arrow[r] \arrow[d] & \wt X^+ \arrow[d] \\
X \arrow[r] & X^+
\end{tikzcd}
\]

则 $\pi_1(X^+)=\pi_1(X)/N$。

对 K 理论，我们作如下选择：$X=B\GL(A)$，$\pi_1(X)=\GL(A)$，$E(A)=N$，且显然 $E(A)\subset\GL(A)$。于是得到 $B\GL(A)^+$，满足
\[ \pi_1(B\GL(A)^+)=\GL(A)/E(A)=K_1(A). \]
我们还有
\[ H_*(B\GL(A),L)\simeq H_*(B\GL(A)^+,L) \]
对所有 $K_1(A)$ 上的模 $L$ 成立；并有如下图表
\[
\begin{tikzcd}
BE(A) \arrow[r] \arrow[d] & \widetilde{B\GL(A)^+} \arrow[d] \\
B\GL(A) \arrow[r] & B\GL(A)^+
\end{tikzcd}
\]

% [原文 p.120]
其中 $\widetilde{B\GL(A)^+}$ 是万有覆盖，而竖直箭头 $BE(A)\to B\GL(A)$ 与 $\widetilde{B\GL(A)^+}\to B\GL(A)^+$ 都是拉回。因此 $BE(A)^+$ 是 $B\GL(A)^+$ 的万有覆盖，从而
\[ \pi_n(B\GL(A)^+)=\pi_n(BE(A)^+),\quad n\ge 2. \]
\end{proof}

对 $n=2$ 有 $\pi_2(B\GL(A)^+)=K_2(A)$。

对 $n=3$ 有 $\pi_3(B\GL(A)^+)=H_3(\mathrm{St}(A))$。

\emph{因此我们有了 Quillen 给出的一般 $K$ 群定义：}

\begin{defn}{1}
\[ K_n(A):=\pi_n(B\GL(A)^+),\quad n\ge 1. \]
\end{defn}

% [Tutor Note: p.119 下图即万有覆盖的拉回图表：$BE(A)\to B\GL(A)$ 与 $\widetilde{B\GL(A)^+}\to B\GL(A)^+$ 均为拉回，故 $BE(A)^+$ 是 $B\GL(A)^+$ 的万有覆盖。原文这几个上标波浪号在 pp.119--122 排得不一致（有时压在 $G$ 上、有时压在整个 $BGLA$ 上），本译本统一写作 $\widetilde{B\GL(A)^+}$。]

% [Tutor Note: p.120 $\pi_n(B\GL(A)^+)=\pi_n(BE(A)^+)$（$n\ge 2$）来自万有覆盖的纤维化 $BE(A)^+\to B\GL(A)^+$（其纤维是离散的 $K_1(A)$，故 $n\ge 2$ 时同伦群相同）。这正是 Quillen 定义一般 $K$ 群的依据：$K_n(A)=\pi_n(B\GL(A)^+)$，$n\ge 1$；$n=1$ 时另由 $\pi_1=K_1(A)$ 给出。]

% [原文 p.121]
\section{第六讲（若干定理与例子）}

下面给出一个不加证明的命题，我们将在下文用它来得到一个基本定义：

\begin{prop}{1} 对映射 $f:X\to Y$（其中 $X,Y$ 是带基点的 CW 复形），下列各项等价：
\begin{enumerate}[label=(\alph*.),leftmargin=2.6em,itemsep=3pt]
\item $f$ 的同伦纤维 $F$ 是无环的，即 $\wt H_*F=0$，其中同伦纤维是指把 $f$ 替换为纤维化后取实际的纤维。
\item $\pi_1(f):\pi_1(X)\to\pi_1(Y)$ 是满射，并且对任意 $\pi_1(Y)$-模 $L$，同调中的诱导映射
\[ f_*:H_*(X,f^*L)\to H_*(Y,L) \]
是同构。
\item 映射 $f:X\to Y$ 的提升 $f':X\times_Y\wt Y\to\wt Y$ 是同调同构（其中 $\wt Y$ 是 $Y$ 的万有覆盖）：
\[
\begin{tikzcd}
X\times_Y\wt Y \arrow[r,"f'"] \arrow[d] & \wt Y \arrow[d] \\
X \arrow[r,"f"'] & Y
\end{tikzcd}
\]
\end{enumerate}
\end{prop}

\begin{defn}{1} 当上述条件成立时，我们称 $f$ 是\emph{无环的}。 \end{defn}

\begin{cor}{1} 无环映射对复合、同伦拉回与同伦推出封闭。 \end{cor}

\begin{thm}{1} 给定完美正规子群 $N\subset\pi_1(X)$，存在唯一的（在同伦等价意义下）无环映射 $f:X\to Y$，使得 $N=\Ker[\pi_1(f)]$（其中 $\pi_1(Y)=\pi_1(X)/N$）。 \end{thm}

% [Tutor Note: p.121 命题 1 原文声明「不加证明」，它是 + 构造的「无环映射」刻画；条件 (a) 中的 $\wt H_*F=0$ 指约化同调为零，即纤维 $F$ 无环（acyclic）。注意 (a.)(b.)(c.) 带括号加句点的编号是原文格式，译文保留。]

% [原文 p.122]
此外，对任意空间 $T$，诱导映射
\[ f^*:[Y,T]\to\{a\in[X,T]\mid\pi_1(a):\pi_1(X)\to\pi_1(T)\ \text{消灭}\ N\} \]
是同构，其中
\[
\begin{tikzcd}
X \arrow[r,"f"] \arrow[d] & Y \arrow[d,equal] \\
T & Y \arrow[l]
\end{tikzcd}
\]

\begin{proof}
关键在于 $Y$ 的构造。设 $\wt X$ 是 $X$ 的覆盖空间，$\pi_1(\wt X)=N$。由于 $f:X\to Y$ 是无环的，我们使用无环推出 $\wt X\to\wt X^+$，即
\[
\begin{tikzcd}
\wt X \arrow[r] \arrow[d] & \wt X^+ \arrow[d] \\
X \arrow[r] & Y
\end{tikzcd}
\]
\end{proof}

如果把它应用到 $E(A)$（它是 $\GL(A)=\pi_1(B\GL(A))$ 的完美正规子群），我们得到：
\[
\begin{tikzcd}
BE(A) \arrow[r,"\tilde f"] \arrow[d] & \widetilde{B\GL(A)^+} \arrow[d] \\
B\GL(A) \arrow[r,"f"'] & B\GL(A)^+
\end{tikzcd}
\]

其中 $f:B\GL(A)\to B\GL(A)^+$ 是满足 $\Ker[\pi_1(f)]=E(A)$ 的唯一无环映射，$\widetilde{B\GL(A)^+}$ 是万有覆盖，且 $\pi_1(B\GL(A)^+)=\GL(A)/E(A)=K_1(A)$。若

% [Tutor Note: p.122 条件 (b) 中「$\pi_1(a)$ 消灭 $N$」的意思是 $N\subset\Ker[\pi_1(a)]$；这个同构正是 + 构造的泛性质。图中右侧竖箭头是恒等映射（双线），表示 $Y$ 同时出现在上下两行。]

% [原文 p.123]
$f$ 是无环的，则 $\tilde f:BE(A)\to\widetilde{B\GL(A)^+}$ 也是无环的。

由上述一切，我们得到：
\begin{itemize}[leftmargin=1.6em,itemsep=2pt]
\item $B\GL(A)^+$ 的万有覆盖是 $BE(A)^+$。
\item $\pi_1(B\GL(A)^+)=K_1(A)$。
\item $\pi_2(B\GL(A)^+)=\pi_2(BE(A)^+)=H_2(BE(A)^+)=H_2(E(A))=K_2(A)$。
\item $\pi_3(B\GL(A)^+)=H_3(\mathrm{St}(A))$。
\item 一般地，$\pi_n(B\GL(A)^+)=K_n(A)$，$n\ge 1$。
\end{itemize}

下面我们来看它与拓扑 K 理论的关系（见第 3 章）。

令
\[ G_\infty=\lim_{k\to\infty}\mathrm{Gr}_k(\C^\infty)=B\GL(\C), \]
其中 $\mathrm{Gr}$ 表示 \emph{Grassmann 流形}，且 $\mathrm{Gr}_k(\C^\infty)=B\GL_k(\C)$。我们假定群 $\GL(\C)$ 带有其自然拓扑。

向量丛的直和 $\xi\oplus\eta$ 对应 $G_\infty$ 上的一个 H-空间结构。注意，虽然 $B\GL(A)$ 不是 H-空间，但 $B\GL(A)^+$ 是。存在连续映射 $\mu:G_\infty\times G_\infty\to G_\infty$（以 $x_0$ 为基点），使得 $\mu(x_0,x)=x$ 且 $\mu(x,x_0)=x_0$。

回忆如下事实：设 $a\in\GL_n(\C)$。映射
\[ a\mapsto\begin{pmatrix} a & 0 \\ 0 & 1 \end{pmatrix} \]
与
\[ a\mapsto\begin{pmatrix} 1 & 0 \\ 0 & a \end{pmatrix} \]
把群 $\GL_n(\C)$ 映入 $\GL_{2n}(\C)$，且它们在自然

% [Tutor Note: p.123 原文 $\mu(x_0,x)=x$ 且 $\mu(x,x_0)=x_0$；按 H-空间的单位元公理，两式都应为 $\mu(x_0,x)=\mu(x,x_0)=x$，原文第二式疑为笔误，正文照原样翻译。]

% [原文 p.124]
拓扑下同伦；但映射
\[ a\mapsto\begin{pmatrix} a & 0 & 0 \\ 0 & 1 & 0 \\ 0 & 0 & 1 \end{pmatrix} \]
与
\[ a\mapsto\begin{pmatrix} 1 & 0 & 0 \\ 0 & a & 0 \\ 0 & 0 & 1 \end{pmatrix} \]
经由 $E(A)$ 中的一个元素是\emph{共轭}的。$E(A)$ 中元素在 $B\GL(A)$ 上的共轭不是平凡的（如果把保基点映射也计入的话），但因为
\[ \pi_1(B\GL(A))=\GL(A), \]
$E(A)$ 中元素在 $B\GL(A)^+$ 上的共轭在同伦意义下是平凡的。

我们知道，同伦的空间在拓扑 K 理论中具有同构的向量丛，但这在代数 K 理论中\emph{不成立}。

一般地，高阶代数 K 理论的计算是困难的。自从 Quillen 定义这些群以来已过去 30 多年，而已知的例子很少。我们给出一些：设 $F[x_1,x_2,\dots,x_n]\sim F$ 是系数取自某个域 $F$ 的多项式环。Grothendieck 证明了
\[ K_0(F[x_1,x_2,\dots,x_n])=K_0F, \]
而 Quillen 证明了多项式环上的每个向量丛都是平凡的。

设 $\F_q$ 是含 $q$ 个元素的有限域。对于这种情形，下列群已被（Quillen）算出：
\[ K_n(\F_q)=\Z,\quad n=0. \]
\[ K_n(\F_q)=\F_q^+\simeq\Z_{q-1},\quad n=1. \]
\[ K_n(\F_q)=0,\quad n=2. \]
\[ K_n(\F_q)=\Z_{q^2-1},\quad n=3. \]
\[ K_n(\F_q)=0,\quad n=4. \]

% [Tutor Note: p.124 原文 $F[x_1,\dots,x_n]\sim F$ 中的 "$\sim F$" 未作解释（大概想表达该多项式环是 $F$ 上的代数），正文照原样保留。又：两族矩阵映射的对比正是 $B\GL(A)^+$ 的 H-空间结构的来源——3×3 情形下两个嵌入经由初等矩阵共轭，而 $E(A)$-共轭在 $B\GL(A)^+$ 上同伦平凡，故两种嵌入给出同一个映射。]

% [原文 p.125]
\[ K_n(\F_q)=\Z_{q^3-1},\quad n=5. \]
\[ \text{\emph{等等}} \]
用更简洁的形式写出，就是：
\[ K_0(\F_q)=\Z, \]
\[ K_{2n}(\F_q)=0,\quad n\in\N^* \]
以及
\[ K_{2n-1}(\F_q)=\Z_{q^n-1},\quad n\in\N^*, \]
其中 $\Z_r$ 表示 $r$ 元循环群。

我们有如下映射：
\[ B\GL(\F_q)\to BU\to BU, \]
其中第一个映射是\emph{Brauer 提升}，$BU=G_\infty$ 是拓扑中的空间，而第二个映射是 $\Psi^q(L)=L^q$。存在从 $B\GL(\F_q)^+$ 到
\[ BU\xra{\Psi^q-1}BU \]
的同伦纤维的一个映射。

我们用到了下述定理：

\begin{thm}{2} 映射 $B\GL(\F_q)^+\to$ 同伦纤维$\bigl(BU\xra{\Psi^q-1}BU\bigr)$ 是同伦等价。 \end{thm}

因为 $B\GL(A)^+$ 是 H-空间，人们知道 $H_*(B\GL(A)^+;\Q)$ 是 Hopf 代数。此外还知道
\[ \pi_*(B\GL(A)^+)\otimes\Q=\mathrm{Prim}[H_*(\GL(A);\Q)], \]
这就是 \emph{Milnor--Moore 定理}。

对 $A=\Z$ 的情形，Borel 证明了：
\[ \dim_\Q[K_n\Z\otimes\Q]=\Z,\quad n=0,5,9. \]
\[ \dim_\Q[K_n\Z\otimes\Q]=0,\quad n=1,2,3,4,6,7,8. \]

% [Tutor Note: p.125 Borel 的公式 $\dim_\Q[K_n\Z\otimes\Q]=\Z$ 有笔误：维数应为 $1$（$K_{4k+1}\Z\otimes\Q\cong\Q$，故 $\dim_\Q=1$）；500 dpi 复核确认原文确实印作 $\Z$，正文照原样保留并在此注明。另 "h-fibre" 是 homotopy fibre 的简写，译文作「同伦纤维」。]

% [Tutor Note: p.125 Brauer 提升是从 $B\GL(\F_q)$ 到 $BU$ 的映射；$\Psi^q$ 是 Adams 运算，对线丛有 $\Psi^q(L)=L^q$。Quillen 的定理 2 说 $B\GL(\F_q)^+$ 恰好是 $\Psi^q-1$ 的同伦纤维，这把有限域上的高阶代数 K 理论完全归约为拓扑 K 理论（$BU$）的计算。]

% [原文 p.126]
更具体地说：

\begin{thm}{3（D. G. Quillen）} 若 $A$ 是代数数域 $\F$（有理数的有限扩张）中的代数整数环，则群 $K_n(A)$ 是有限生成的。 \end{thm}

Borel 用这个结果计算了 $K_n(A)$ 与 $K_n(\F)$ 模去挠的部分。例如，对于整数，Borel 证明了（模挠）
\[ K_n(\Z)=0 \]
对正的 $n$ 成立，除非 $n=4k+1$ 且 $k$ 也是正的；并且（模挠）
\[ K_{4k+1}(\Z)=\Z \]
对正的 $k$ 成立。$K_{2n+1}(\Z)$ 的挠子群与有限群 $K_{4k+2}(\Z)$ 的阶最近已被确定，但后面这些群是否为循环群、以及群 $K_{4k}(\Z)$ 是否为零，都依赖于所谓 \emph{Vandiver 猜想}（关于分圆整数的类群）。

% [Tutor Note: p.126 原文 "whether the later groups are cyclic" 中的 "later" 应为 "latter"（指 $K_{4k+2}(\Z)$ 这批有限群），正文照译。Vandiver 猜想大致是说：奇素数 $p$ 对应的分圆域 $\Q(\zeta_p)$ 的实子域的类数不被 $p$ 整除。]

% [Tutor Note: p.126 原文列举 $K_n(\Z)$ 的非零（模挠）位置为 $n=0,5,9$，这与 Borel 定理一致：$K_0(\Z)=\Z$，且 $K_{4k+1}(\Z)$（$k\ge1$）贡献 $\Z$；$K_1(\Z)=\Z/2$、$K_{4k+2}(\Z)$ 等只有挠，故 $\dim_\Q=0$。据此可核对上面两行 $\dim_\Q$ 的列表。]
